% PHY 210 Syllabus, Smith College, Fall 2026 -- Michael Lerner
%
% Rewritten 2026-08-17 from the Earlham 2024 syllabus (see git history for
% the original). Sources: the F2026 calendar (make_fall2026_calendar.py, the
% source of truth for dates and the 1000-point scheme), the S26 Smith
% syllabus (boilerplate: Honor Code, ARC, land acknowledgment, resources),
% Will Raven's F2025 AI statement (adapted, credited), Gary Felder's PCCI
% policy language (adapted), and Michael's own Earlham prose (course goals,
% communication, workload, sensitive issues, diversity, late passes).
%
% Red [TODO: ...] markers flag facts to fill in before distribution.
% Build: pdflatex PHY210Syllabus.tex   (from this directory)
\documentclass[12pt]{article}
%%%%%%%%%%%%%%%%%%%%%%%%%%%%%%%%%%%%%%%%%%%%%%%%%%%%%%%%%%%%%
\input{../macros}
%%%%%%%%%%%%%%%%%%%%%%%%%%%%%%%%%%%%%%%%%%%%%%%%%%%%%%%%%%%%%
\usepackage[table]{xcolor}
\usepackage{multirow}
\usepackage{hyperref}
\hypersetup{
    colorlinks=true,
    linkcolor=blue,
    filecolor=magenta,
    urlcolor=blue,
  }
\usepackage{float}
\usepackage[super]{nth}
\usepackage{booktabs}
\restylefloat{table}

\newcommand{\todonote}[1]{{\color{red} \textbf{[TODO: #1]}}}

\pagestyle{empty}

\renewcommand{\thefootnote}{\fnsymbol{footnote}}
\begin{document}

\begin{center}
{\bf PHY 210: Mathematical Methods of Physical Sciences and Engineering\\
Smith College, Fall 2026}
\end{center}

\setlength{\unitlength}{1in}

\hrule

\vskip.15in
\noindent\textbf{Instructor:} Michael Lerner, \todonote{office},
Email: \href{mailto:mglerner@smith.edu}{mglerner@smith.edu}

\begin{description}
  \item[Office Hours:] \todonote{days/times, set after the week-1
        scheduling poll}. I also have an open-door policy, and you're
        encouraged to stop in to ask questions whenever my door is
        open. That's most of the time.
  \item[When and where:] Class meetings are MWF,
        \todonote{time and room}. First meeting Wednesday, September 9;
        last meeting Monday, December 14.
  \item[Required materials:] Felder \& Felder, \textit{Mathematical
        Methods in Engineering and Physics} (Wiley). Some sections and
        all answers to odd problems are free online at
        \url{https://www.felderbooks.com/mathmethods/}. We will also use
        Python (via Jupyter notebooks on Smith's server at
        \url{https://posit.smith.edu} -- nothing to install or buy).
        Course announcements, documents, and assignments live on Moodle.
  \item[Prerequisites:] MTH 212 (multivariable calculus) and PHY 117 or
        PHY 119, or permission of the instructor. \textbf{Differential
        equations are NOT a prerequisite} -- this course is where you
        learn them.
  \item[Course Goals:]
        One of the great strengths of a modern
        physics background is that it allows you to address problems in an
        enormous range of fields, from physics itself to economics, biology,
        ecology, computer science, complex systems, physical chemistry,
        geology, and engineering. The primary objective of this course is to
        provide a systematic introduction to mathematical techniques that will
        serve you both in future physics courses and in modeling interesting
        problems in your domain of choice. A significant secondary objective
        is to develop a level of mathematical sophistication that will allow
        you to confidently and competently explore further material on your
        own.

        Concretely, by the end of this course you will be able to:
        \ \\\vspace{-.3in}
        \begin{enumerate}
          \item Describe a physical scenario mathematically (most often
                as a differential equation), solve the resulting model,
                verify your solution, and interpret it back in the real
                world -- including honestly assessing where the model
                breaks down.
          \item Apply the core toolboxes of mathematical physics:
                ordinary differential equations, Laplace transforms,
                complex numbers, series approximations, linear algebra
                through eigenvectors and normal modes, multidimensional
                integrals, vector calculus, Fourier series, and an
                introduction to partial differential equations.
          \item Check your work the way a scientist does: verify units,
                plug solutions back into the original equations, and
                test limiting cases.
          \item Use Python to perform symbolic calculations, and to
                visualize and explore mathematical models.
        \end{enumerate}
\end{description}

\newpage

\noindent\textbf{Course Calendar}

\vskip.1in
\noindent The complete day-by-day schedule (topics, reading due,
homework due) is posted on Moodle and will be kept up to date there.
The shape of the semester:

\setlength{\arrayrulewidth}{.4mm}
\setlength{\tabcolsep}{8pt}
{\rowcolors{2}{green!80!yellow!50}{green!70!yellow!40}
  \begin{table}[H]
    \footnotesize
\begin{tabular}{l|p{0.52\linewidth}|l}
\textbf{Weeks} & \textbf{Topics} & \textbf{Felder \& Felder} \\ \hline
1-3   & Ordinary differential equations: setting them up from physical
        situations, arbitrary constants, separation of variables, guess
        and check, superposition & Ch.\ 1 \\
3-4   & Heaviside, Dirac delta, and the Laplace transform (with Python)
      & 10.2, 10.10-10.11 \\
5     & Complex numbers and Euler's formula & Ch.\ 3 \\
6-7   & Linear approximations, Maclaurin and Taylor series & Ch.\ 2 \\
7-9   & Linear algebra: matrices as transformations, eigenvalues and
        eigenvectors, coupled oscillators and normal modes & Ch.\ 6 \\
10-11 & Integrals in two and more dimensions; cylindrical and spherical
        coordinates; line and surface integrals & Ch.\ 5 \\
11-13 & Vector calculus: gradient, divergence, curl -- \textbf{defined
        geometrically, following Feynman} -- then the divergence theorem,
        Stokes' theorem, and conservative fields & Feynman II 2-3; Ch.\ 8 \\
14    & Fourier series & Ch.\ 9 \\
15    & Partial differential equations: the heat equation, separation of
        variables, normal modes of the wave equation; Fourier transforms
        as the closing vista & Ch.\ 11; 9.6 \\
\end{tabular}
\end{table}
}

\noindent\textbf{Key dates:}
\begin{itemize}
  \item Math pretest due Friday, September 18 (graded for completion).
  \item Quizzes, in class, full period: Friday September 25 (Ch.\ 1);
        Friday October 23 (Ch.\ 10, 3, 2); Friday November 20
        (Ch.\ 6, 5).
  \item Non-Newtonian Scientist assignment due Monday, October 26.
  \item Friday October 9 is a flex day (it absorbs Mountain Day, snow,
        or wherever we're behind).
  \item No class: Monday October 12 (autumn recess), Wednesday November
        25 and Friday November 27 (Thanksgiving).
  \item Final exam during the registrar's exam period, December 19-22.
\end{itemize}

\newpage

\begin{description}
  \item[Recommended Textbooks] \hfill \\
        \textbf{The Feynman Lectures on Physics, Volume II} -- Chapters
        2 and 3 are \emph{assigned reading} for our vector-calculus
        unit: we will meet divergence and curl through Feynman's
        physical, geometric definitions before we see the formulas in
        Felder. This part of Feynman is extremely conversational and
        readable, and you can read the whole thing free, with beautiful
        typesetting, at \url{https://www.feynmanlectures.caltech.edu/}. \\
        \textbf{Boas, Mathematical Methods in the Physical Sciences} --
        the classic alternative presentation. If Felder's treatment of a
        topic isn't clicking, try Boas's; I taught from her book for a
        decade and am happy to point you at the parallel section. \\
        \textbf{Schey, div grad curl and all that} -- an extremely
        conversational introduction to (and refresher on) vector
        calculus. \\
        \textbf{Boyce and DiPrima, Elementary Differential Equations and
        Boundary Value Problems} -- a thorough reference if you want
        more depth on differential equations than we cover. \\
        \textbf{Arfken and Weber, Mathematical Methods for Physicists}
        -- an encyclopedic reference; too dry to learn from in this
        course, but useful later in your career.

  \item[Communication:] My preferred mode of communication is email. I
        will check my email multiple times per day during the work
        week. If you have sent me an email and have not heard back
        from me within 24 hours (excluding weekends), you should email
        me again or drop by my office to ask your question in
        person. I cannot guarantee that I will regularly check my
        email on evenings or on weekends.

        I expect that you will check your Smith email at least once per
        day. That is, 24 hours after sending an email to the class, it
        is my expectation that all of you will have read that
        email. Emailed assignments, due dates, policy changes, and
        other course information carry the same weight as information
        contained in the syllabus. Information relevant to the course
        will also be posted to Moodle, and you are responsible for
        checking Moodle with the same frequency.

  \item[Workload Expectations:] This course should take about 12 hours
        of your time per week: about 4 in class, and about 8 outside it
        on reading, pre-class check-ins, homework problems, and review.
        \textbf{If you find yourself spending significantly more (or
        less) than this, please let me know ASAP! That's an
        instructor-error that I can easily fix.}

  \item[Contacting me about sensitive issues:] What if you have
        something you want to talk about, but don't feel comfortable
        doing so in a public (or non-anonymous) space? This often
        comes up with issues of racism, sexism, and other kinds of
        discrimination. It can come up with issues that relate to my
        behavior, the behavior of your classmates, or the world in
        general. Please feel {\it encouraged} to use either the
        anonymous feedback form linked on Moodle, or to leave a written
        note under my door. If you want to write a note, you may want
        to type it up and print it out so that your handwriting is
        further anonymized. Of course, if you feel comfortable and safe
        doing so, you are also encouraged to talk publicly or
        non-anonymously about these things! If you are not comfortable
        with any of these options, the Resources section at the end of
        this syllabus lists offices that can help, including
        confidential ones.

  \item[Tutoring and study help:] The Spinelli Center supports
        quantitative courses across the college, with both drop-in
        hours and one-on-one appointments, plus math review workshops.
        \todonote{course tutor / Learning Assistant names and tutoring
        hours, once assigned}.

  \item[Diversity Statement:] I can hardly overstate how much inclusion,
        equity and diversity considerations have shaped my entire thinking
        around teaching and practicing physics. It affects everything from the
        way that physics is framed historically to the way that it is
        typically taught in your classrooms. My main goal is to create an
        environment that tells all of us that we can do physics; that invites
        people into physics who have never been invited, or who have been
        explicitly excluded. In this class, the ``Non-Newtonian
        Scientist'' assignment (due Monday, October 26) engages directly
        with who does physics and how the standard physics canon got
        assembled; the full prompt will be posted on Moodle.

        Physics cannot be divorced from the world in which we study it. I
        expect that we will recognize that in this class, and I expect
        that we will create an inclusive and welcoming
        environment. More specifically, I expect that we will strive
        \textit{not} to create an environment that excludes
        people. This is hard work. We will get things wrong. You will
        get things wrong. \textit{I} will get things wrong. Especially
        when I get things wrong, I encourage you to use any of the
        channels above -- direct comments, the anonymous form, notes
        under my door -- to tell me.

        \newpage

  \item[Grading:] Everything in this course is worth points, and the
        course total is exactly 1000 points. Within each category, each
        assignment is weighted equally.

\setlength{\arrayrulewidth}{1mm}
\setlength{\tabcolsep}{8pt}
{\rowcolors{2}{green!80!yellow!50}{green!70!yellow!40}
\begin{table}[H]
\label{tab:overall-grade-weights-table}
\begin{tabular}{lllll}
\hline
\textbf{Category} & \textbf{Number} & \textbf{Drop} & \textbf{Points each} & \textbf{Total} \\ \hline
\textbf{Participation \& pre-class check-ins (PCCIs)} & 39 & 4 & 1 & 35 \\
\textbf{Weekly homework (WHW)} & 13 & 1 & 25 & 300 \\
\textbf{Non-Newtonian Scientist} & 1 & 0 & 25 & 25 \\
\textbf{Quizzes (in class)} & 3 & 0 & 150 & 450 \\
\textbf{Final exam} & 1 & 0 & 190 & 190 \\ \hline
\textbf{Course total} & & & & \textbf{1000} \\ \hline
\end{tabular}
\end{table}
}

  \item[Pre-Class Check-Ins (PCCIs):] Most class days there will be a
        short PCCI due at the start of class, on paper. Many of these
        are the textbook's ``Discovery Exercises,'' designed to
        introduce you to a topic \emph{before} we cover it; others are
        short problems on material we've just covered. \textbf{You will
        receive full credit for a good-faith effort that includes a
        clear description of anything you are stuck on.} An incorrect
        or incomplete answer \emph{without} such an explanation won't
        receive full credit. For example, if the question were to
        differentiate $f(g(x))$: ``I don't know'' earns nothing; the
        correct answer earns full credit; and so does ``I think it's
        $f'(g(x))$, except I know the chain rule says to do something
        different because it's $f(g(x))$ and not just $x$ -- I'm not
        clear on what, though.'' A PCCI should take you less than 15
        minutes; if one doesn't, please tell me. Turning in the day's
        PCCI is also how attendance is recorded, and your four lowest
        participation days are dropped, so occasional absences and
        off days cost you nothing.

  \item[Weekly Homework (WHW):] This is basically a course in
        problem-solving techniques. The most common and the most damaging
        mistake a student can make in such a course is to yield to the
        temptation to try and master the material by reading the text and
        listening to the lectures while working a minimum of problems. While
        reading and listening can give you valuable new ideas, problem-solving
        skills are only efficiently developed by solving as many problems as
        possible. For most students the understanding and knowledge gained
        from this course will be in direct proportion to the number of
        problems which they successfully complete.

        Each week I will post a WHW problem list, organized in three
        tiers: \textbf{Warm-up} (start here), \textbf{Essentials} (the
        core -- if you can do these, you're on track), and
        \textbf{Depth} (for when you want more). You are not expected
        to do every problem; do whichever problems you need practice
        on. If a tier starts to feel easy, that's a good sign that
        you're comfortable with it. Answers to odd problems are in the
        book's Appendix M
        (\url{https://www.felderbooks.com/mathmethods/}).

        Each Friday by the end of the day, you will turn in (on Moodle):
        (1) a photograph of your written work on \textbf{at least one
        problem} from that week's list, with a sentence or two about
        where you got stuck or what you'd like feedback on, and (2) a
        short reflection form. WHWs are graded on the same good-faith
        standard as PCCIs: thoughtful engagement earns full credit.
        \todonote{confirm final WHW turn-in structure before
        distribution}

  \item[Quizzes:] Three quizzes, in class, taking the full period, on
        the Fridays listed above. Sample quizzes with solutions will be
        posted beforehand, and quiz questions will be closely aligned
        with the in-class problems and the WHW lists. Quizzes are
        closed-book; you may bring \textbf{one 8.5$\times$11 sheet of
        handwritten notes, both sides}. \todonote{calculator policy --
        recommend: a basic calculator is fine; no phones, laptops, or
        internet-connected devices} Each quiz covers material from
        earlier weeks, never the week the quiz is given, and I will
        announce the exact section coverage in class about a week
        ahead. If you need accommodations, contact me well before the
        quiz.

  \item[Final exam (with redemption problems):] The final has two
        parts. The \emph{required} part covers the last third of the
        course (vector calculus, Fourier series, and PDEs). The rest
        consists of \emph{optional redemption problems} drawn from the
        three quizzes: if you do better on a redemption problem than
        you did on the corresponding quiz problem, the new score
        replaces the old one. There is no penalty for attempting a
        redemption problem -- you cannot lose points. Do the required
        problems first; then do as many redemption problems as you
        want and have time for.

  \item[Partial credit, and how to get it:] Several of the problems in
        this class are quite challenging. For the harder problems, my
        goal is to have you make the strongest possible effort toward
        \textbf{understanding} the solution. Thus, if you cannot fully
        solve a problem (on a quiz or the final), say whatever you can
        about the way in which a solution would proceed from where you
        stop; say whatever you can about the qualitative behavior of a
        solution; say whatever you can about the physical meaning of
        the solution. And always, always: \textbf{check your units},
        and \textbf{plug your solution back into the original
        equation}. You are never done solving a differential equation
        until you have verified that your solution works.

  \item[Late policy:] It is extremely hard to catch up on late work
        (usually things are late because you are busy, and people
        don't tend to get less busy!). That is why every category has
        built-in drops: four participation days, one full WHW. My
        strong suggestion, when life happens, is to use the drops and
        prioritize getting your schedule back under control.

        Beyond the drops: all students start the term with \textbf{3
        free late passes}. You can spend a late pass to extend a WHW
        (or the Non-Newtonian assignment) to the following class period
        with no penalty. Late passes are not automatic -- email me
        before the deadline to say you're using one, and once declared
        it is spent whether or not you turn the work in by the extended
        date. Late passes cannot be applied to PCCIs (their whole value
        is priming you for that day's class -- the drops cover them),
        to quizzes, or to the final. After your free passes are spent,
        further extensions cost 10\%, then 20\%, and so on, of that
        assignment's points.

        \newpage

  \item[Grade scale:] Grades are not curved. Final grades follow the
        standard scale:

\setlength{\arrayrulewidth}{1mm}
\setlength{\tabcolsep}{8pt}
{\rowcolors{2}{green!80!yellow!50}{green!70!yellow!40}
\begin{table}[H]
\label{tab:overall-grade-cutoffs-table}
\begin{tabular}{llll}
\hline
 \textbf{Grade} & \textbf{Range} & \textbf{Grade} & \textbf{Range} \\ \hline
A  & 93-100\% & C+ & 77-79.9\% \\
A- & 90-92.9\% & C  & 73-76.9\% \\
B+ & 87-89.9\% & C- & 70-72.9\% \\
B  & 83-86.9\% & D+ & 67-69.9\% \\
B- & 80-82.9\% & D  & 63-66.9\% \\
   &           & D- & 60-62.9\% \\
   &           & E  & 0-59.9\% \\
\end{tabular}
\end{table}
}

  \item[Use of Artificial Intelligence:] AI is a powerful tool, and it
        offers real potential to enhance your learning in this course.
        Used improperly, it can genuinely hinder your development as a
        physicist. You are welcome to use AI in this class, provided
        you always adhere to the Honor Code: \textbf{using AI to solve
        homework, PCCI, quiz, or exam problems is a violation of the
        Honor Code and will be treated as such}, and any AI use in
        work you submit must be cited, like any other source.

        Two things are worth saying plainly. First, AI models make
        subtle errors, hallucinate, and hand you technically-correct
        solutions that teach you nothing; your job is to verify,
        question, and ultimately understand the physics. Second --
        and this matters in \emph{this} course specifically -- your
        PCCIs and WHWs are graded on good-faith effort, not
        correctness. There are no points available for a polished
        answer. The only thing being graded is the record of
        \emph{your} thinking, including where you got stuck; having an
        AI polish that record destroys the one thing it's for. Where
        AI genuinely helps: asking for alternative explanations of a
        concept, checking your algebra \emph{after} you've done it,
        and exploring beyond the course. [Adapted, with thanks, from
        Will Raven.]

  \item[Recording:] \todonote{finalize recording policy -- planned:
        class projections are captured as PDFs and posted after each
        class; some or all classes may be video-recorded and posted for
        enrolled students}

  \item[Academic Integrity and Honesty / Honor Code:] Students and
        faculty at Smith are part of an academic community defined by
        its commitment to scholarship, which depends on scrupulous and
        attentive acknowledgement of all sources of information and
        honest and respectful use of college resources. Smith College
        expects all students to be honest and committed to the
        principles of academic and intellectual integrity in their
        preparation and submission of course work and examinations. All
        submitted work of any kind must be the original work of the
        student, who must cite all the sources used in its preparation
        (including, for example, any AI tools). Students voted to
        establish the academic honor system in 1944. The basis of the
        Academic Honor Code is articulated in Article X of the SGA
        Constitution and Article VII of the SGA Bylaws.

        In this class specifically: you are encouraged to work together
        on the WHW problems and to study together for quizzes, but
        PCCIs should be attempted individually first, everything you
        turn in must be your own work, and quizzes and the final must
        be done individually.

  \item[Academic Accommodations:] If you need classroom accommodations,
        please contact me and include your letter from the
        Accessibility Resource Center (ARC,
        \href{mailto:arc@smith.edu}{arc@smith.edu}). I promise to be
        flexible in meeting any documented need for accommodations in a
        way that works well for you. If you would like some
        accommodation but do not have an ARC letter, let me know and we
        can discuss your concerns: with or without a letter, I want to
        provide you the support you need to succeed in this course.

  \item[Land Acknowledgment:] We acknowledge that we are on Indigenous
        land: the territory of the Nipmuc people. We are grateful for
        the opportunity to live, learn, and grow on this sacred land,
        and extend our respect to citizens of this nation who live here
        today, and to their ancestors who have lived here for hundreds
        of generations. We recognize the repeated violations of
        sovereignty, territory, and water perpetrated by invaders who
        have impacted the original inhabitants of this land for 400
        years. We know this acknowledgement is insufficient, and does
        not undo the harm that has been done and continues to be
        perpetrated now against Indigenous people and their land and
        water.

  \item[Course feedback:] During the last week of classes,
        approximately 15 minutes at the beginning of one class will be
        set aside for you to complete the course feedback
        questionnaire.

  \item[Resources:] There are many resources at Smith to help you
        succeed in this course and in your college career. Here are
        some:
        \begin{itemize}
          \item Your instructor! I am here to help you. That's my job,
                so don't think it's bothering me to have you reach out
                with questions or concerns.
          \item The Spinelli Center -- open hours, one-on-one
                appointments, and math review workshops.
          \item Workshops for time management and managing stress
                (Tutoring \& Learning Specialist Services).
          \item Writing help from the Jacobson Center.
          \item Crisis resources, counseling resources, and wellness
                services (Schacht Center).
          \item Resources for gender identity and expression.
          \item Where to report sexual misconduct and other forms of
                discrimination. Please note that your instructors are
                Responsible Reporters.
        \end{itemize}
\end{description}

\end{document}
